%! Exponential Function Manipulation — Unit 2, Lesson 2.4 — Warm-Up (Answer Key)
\documentclass[10pt]{article}
\usepackage{precalculus-article}
\usepackage{precalculus-key}   % pulls in -boxes; do NOT also load -boxes
\newcommand{\TallMath}[1]{$\displaystyle #1\rule[-1.4em]{0pt}{3.2em}$}

\begin{document}

\pageheader{Unit 2, Lesson 2.4}{Warm-Up --- Answer Key}
\namedateperiod

\vspace{0.14in}

\begin{enumerate}[itemsep=16pt]

\item For $f(x) = 5 \cdot (0.4)^x$, identify the values of $a$ and $b$, then classify the
      function as \textbf{exponential growth} or \textbf{exponential decay}.

      \vspace{0.06in}
      $a = $ \ans{5} \quad $b = $ \ans{0.4} \quad Circle one: \textbf{Growth} \;/\; \ans{Decay}

      \begin{teachernote}
        $b = 0.4 < 1$ so the function decays. $a = 5$ is the initial value ($f(0) = 5$).
      \end{teachernote}

\item Evaluate $f(2)$ for $f(x) = 3 \cdot 2^x$. Show your work.

      \vspace{0.06in}
      $f(2) = 3 \cdot 2^2 = 3 \cdot 4 = $ \ans{12}

\item The table below shows values of a function. Write an equation for $f(x)$ in the form
      $ab^x$.

      \vspace{0.08in}
      \begin{center}
      \renewcommand{\arraystretch}{1.4}
      \begin{tabular}{|c|c|}
        \hline
        $x$ & $f(x)$ \\
        \hline
        0 & 12 \\
        1 & 4 \\
        2 & $\dfrac{4}{3}$ \\
        3 & $\dfrac{4}{9}$ \\
        \hline
      \end{tabular}
      \end{center}

      \vspace{0.06in}
      $a = $ \ans{12} \quad $b = $ \ans{$\dfrac{1}{3}$} \quad $f(x) = $ \ans{$12\cdot\left(\dfrac{1}{3}\right)^x$}

      \begin{teachernote}
        Common ratio: $4 \div 12 = \tfrac{1}{3}$. Initial value from table: $f(0) = 12$.
      \end{teachernote}

\end{enumerate}

\end{document}
